\def\stat{l2-7}





\def\tit{НОВЫЕ ОСОБЕННОСТИ САМОРАСПРОСТРАНЯЮЩИХСЯ 


ВРЕДОНОСНЫХ ПРОГРАММ}








\def\titkol{Новые особенности самораспространяющихся 


вредоносных программ} 





\def\autkol{М.\,В.~Левыкин}





\def\aut{М.\,В.~Левыкин$^1$}





\titel{\tit}{\aut}{\autkol}{\titkol}





%{\renewcommand{\thefootnote}{\fnsymbol{footnote}}\footnotetext[1]


%{Работа выполнена при поддержке РФФИ, грант 10-01-00480.}}





\renewcommand{\thefootnote}{\arabic{footnote}}


\footnotetext[1]{Институт проблем информатики Российской академии наук, de\_shiko@yahoo.com}


      


      \Abst{Проведенные исследования вредоносных программ, таких как червь 


{Stuxnet} и буткит TDL/TDSS, позволили выявить следующую особенность: 


современные вредоносные программы представляют собой многоагентные системы, 


имеющие сложную структуру на различных уровнях операционной системы (ОС). Такие 


многоагентные системы могут состоять из следующих агентов: скрытия, коммуникации, 


функционирования, распространения и~т.\,п. Интерфейс взаимодействия между агентами 


таких систем представляет особый интерес.}


      


      \KW{вредоносная программа; сетевой червь; {Stuxnet}; TDS/TDLL}


      


                        \vskip 12pt plus 9pt minus 6pt








      \thispagestyle{headings}





      


            \label{st\stat}


            


            \vspace*{-9pt}





\section{Введение}





     Целью настоящей работы является выявление новых особенностей в 


архитектуре современных самораспространяющихся вредоносных программ. 


Для достижения поставленной цели в работе решаются следующие задачи: 


дается обзор характерных особенностей архитектуры традиционных 


самораспространяющихся вредоносных программ, проводится исследование 


современных самораспространяющихся вредоносных программ и особенностей 


их архитектуры. 





\section{Традиционная архитектура самораспространяющихся 


вредоносных программ в~операционной~системе~Windows}


     


     Традиционный сетевой червь представляет собой монолитную 


программу, которая, эксплуатируя уязвимость в используемом программном 


обеспечении (ПО) пользователя или методы социальной инженерии, 


производит несанкционированное копирование самой себя в операционную 


систему пользователя для выполнения вредоносного воздействия. При этом 


процесс распространения не контролируется вредоносной программой, т.\,е.\ 


происходит хаотично. Традиционно червь содержит в себе примитивную 


функциональную часть, например обращение на определенный сетевой ресурс 


для реализации {DDoS}-атаки. Такие самораспространяющиеся 


программы являются пользовательскими приложениями, так как способы их 


распространения зависят от эксплуатируемых уязвимостей в ПО. 


Большинство самых знаменитых сетевых червей до 2010~г.\ были 


реализованы на пользовательском уровне в ОС в виде 


единых программных модулей~[1].


     


     Пути проникновения самораспространяющейся программы зависят от 


приложения. Можно выделить некоторые наиболее популярные 


самораспространяющиеся вредоносные программы в зависимости от путей 


проникновения~[2]:


     \begin{itemize}


\item почтовые черви, распространяющиеся в формате сообщений 


электронной почты, например в виде ссылок на зловредные сетевые 


ресурсы, или использующие уязвимости в почтовом клиенте пользователя;


\item интернет-черви, использующие для распространения протоколы 


глобальной сети Интернет. Этот тип червей преимущественно 


распространяется с использованием неправильной обработки некоторыми 


приложениями базовых пакетов стека протоколов TCP/IP. 


Например, известное семейство червей conficker~[3] использует 


уязвимость в службе удаленного вызова процедур (RPC) ОС 


Windows;


\item {LAN}-черви, распространяющиеся по протоколам локальных 


вычислительных сетей, например в результате эксплуатирования 


уязвимостей в реализации протокола {NetBIOS}.


\end{itemize}





\section{Выявленные особенности архитектуры некоторых новых 


самораспространяющихся вредоносных программ в~операционной


системе Windows}


     


     В данной работе рассмотрены две современные самораспространяющиеся 


вредоносные программы: червь {Stuxnet} и буткит {TDL/TDSS}.





\subsection{Червь Stuxnet}


     


     Червь {Stuxnet}~--- это сетевой червь нового поколения. 


     Во-пер\-вых, потому что в нем, в отличие от традиционных 


самораспространяющихся программ, используется сразу три способа 


распространения. Во-вто\-рых, потому что он представляет собой комплексную 


систему, направленную на физическое разрушение прог\-рам\-мно-аппаратных 


средств. Данный червь состоит из следующих агентов: скрытия~--- это драйвер 


уровня ядра; распространения~--- агенты, использующие уязвимости в ОС 


Windows; функциональных агентов в виде библиотек для работы с 


Siemens-контроллером, использующимся наиболее часто при работе с ядерным 


реактором. Модуль уровня ядра заменяет оригинальные библиотеки для работы 


с указанным контроллером на свои собственные. Подменные библиотеки 


устанавливают функции перехвата интерфейса работы с контроллером. Затем 


данные функциональные агенты вводят программно-аппаратные комплексы в 


резонирующее состояние, что приводит к их разрушению~[4]. 





\subsection{Буткит TDL/TDSS}





     Другим примером современной самораспространяющейся вредоносной 


сис\-те\-мы является буткит {TDL/TDSS}. Данная многоагентная система 


пошла еще дальше в своем развитии, чем червь Stuxnet. Во-пер\-вых, в 


качестве технологии скрытия в ней используется буткит. Этот буткит с 


помощью метода рекурсивного захвата системных функций осуществляет 


скрытие данной многоагентной системы. Система имеет в своем составе также 


агенты <<полезной нагрузки>>. Эти агенты, выполненные в виде динамических 


библиотек, внедряются в определенный пользовательский процесс. 





\subsection{Архитектурные особенности современных 


самораспространяющихся программ}


     


     Современные вредоносные самораспространяющиеся программы, 


рас\-смот\-рен\-ные в данной работе, обладают новым специфическим общим 


признаком. Они построены по принципу многоагентных систем. Указанные 


вредоносные программы состоят из двух и более групп агентов, реализованных 


на разных уровнях доступа к ОС. Агенты скрытия 


реализованы на уровне ядра ОС, функциональные агенты, или <<агенты 


полезной нагрузки>>, реализованы на уровне пользователя ОС. Каждый из 


агентов, в свою очередь, может состоять из составных частей. При этом для 


скрытия факта загрузки агентов функциональной части используется 


специальный метод. Этот метод позволяет скрыть от средств защиты загрузку 


исполняемых модулей многоагентной системы. Под скрытием в данном случае 


понимается легализация агента, т.\,е.\ формирование условий, при которых 


средства защиты принимают деятельность агента злоумышленника за 


легальный процесс или программу. Примером метода легализации может 


служить изменение системных файлов и их контрольных сумм в ядре с 


внесением программных закладок в эти файлы.





\section{Заключение}





     В данной работе на примерах современных самораспространяющихся 


вредоносных программ~--- червя {Stuxnet} и буткита {TDL/TDSS}~--- 


выявлены новые особенности построения самораспространяющихся 


вредоносных программ. 


     


     На основании проведенных исследований можно сделать следующий 


вывод: современные самораспространяющиеся вредоносные программы 


представляют собой многоагентные системы, состоящие из функциональных 


агентов и агентов скрытия; при этом для скрытия функциональных агентов 


используется метод легализации, основанный на внедрении исполняемого кода 


из агента скрытия режима ядра в пользовательский процесс.


    


{\small\frenchspacing


{%\baselineskip=10.8pt


\addcontentsline{toc}{section}{Литература}


\begin{thebibliography}{9}





\bibitem{1-2-7}


Самые знаменитые компьютерные вирусы и черви. {\sf 


http://www.hardnsoft.ru/ ?trID=1297\&artID=14743}.





\bibitem{2-2-7}


Червь. {\sf 


http://www.tadviser.ru/index.php/\%D0\%A1\%D1\%82\%D0\%B0\%D1\%82\% D1\%8C\%D1\%8F:\%D0\%A7\%D0\%B5\%D1\%80\%D0\%B2\%D1\%8C}.





\bibitem{3-2-7}


\Au{Conficker~A.}


2008. {\sf http://www.securitylab.ru/virus/363745.php}.


 


 \label{end\stat}





\bibitem{4-2-7}


F-secure. Stuxnet questions and answers.


{\sf http://www.f-secure.com/weblog/archives/ 00002066.html}.


\end{thebibliography}


}


}    


